\documentclass[12pt]{article}
\usepackage{fullpage}
\usepackage{amsmath,amssymb,amsthm}

\newtheorem{lemma}{Lemma}
\newtheorem{proposition}{Proposition}
\newcommand{\R}{\mathbb R}
\newcommand{\E}{\mathbb E}
\newcommand{\1}{\mathbf 1}
\newcommand{\Var}{\operatorname{Var}}
\newcommand{\TV}{\operatorname{TV}}
\newcommand{\Law}{\operatorname{Law}}
\newcommand{\calF}{\mathcal F}
\newcommand{\Ent}{\operatorname{Ent}}
\newcommand{\Lip}{\operatorname{Lip}}
\newcommand{\calE}{\mathcal E}

\begin{document}
\title{Submission A solution to Problem 5}
\date{}
\maketitle

\begin{center}
{\bf Uniqueness of the invariant measure}
\end{center}

\section*{Problem statement and interpretation}

Put \(E=C_0([0,1])\), \(H=L^2(0,1)\), and let \(S_t\) be the Dirichlet heat
semigroup generated by \(\frac12\Delta\).  I use the Dirichlet version of the
Athreya--Butkovsky--L\^e--Mytnik approximation assumption in the following
form.  For each deterministic initial condition \(x\in E\), if \(b_n=G_{1/n}\),
then the regularized equation has the mild solution
\[
  u_t^n=S_t x+V_t+K_t^n,\qquad
  K_t^n=\int_0^tS_{t-r}b_n(u_r^n)\,dr,\qquad
  V_t=\int_0^tS_{t-r}\,dW_r ,
\]
and, for every \(t>0\) and every \(\Phi\in C_b(E)\),
\[
        P_t^n\Phi(x):=\E\Phi(u_t^n)\longrightarrow P_t\Phi(x).
        \tag{1}
\]
Thus I interpret the stated weak convergence of \(u^n(t)\) to \(u(t)\) as
holding for all deterministic starting points; this is the form needed for a
statement about invariant measures of the whole Markov semigroup.

\begin{lemma}[Uniform estimate on the regularized drift]\label{lem:Hestimate}
For each \(T<\infty\) there is \(C_T<\infty\), independent of \(n\) and
\(x\in E\), such that
\[
 \|K_t^n-S_{t-s}K_s^n\|_{L^2(\Omega;H)}
      \le C_T |t-s|^{3/4},\qquad 0\le s<t\le T .             \tag{2}
\]
\end{lemma}

\begin{proof}
Let \(p_t\) be the Dirichlet kernel and \(d(y)=y\wedge(1-y)\).  For
\(0<a\le 2T\), the reflection formula for the interval and comparison with
the two half-lines at the nearer boundary give
\[
p_a(z,y)\le Ca^{-1/2}e^{-|z-y|^2/(Ca)},\qquad
p_a(y,y)\asymp a^{-1/2}\bigl(1\wedge d(y)^2/a\bigr).          \tag{3}
\]
For example, if \(d(y)\le\sqrt a\), the half-line expression
\((2\pi a)^{-1/2}(1-e^{-2d(y)^2/a})\) is comparable to \(d(y)^2a^{-3/2}\);
if \(d(y)\ge\sqrt a\), the reflected terms are exponentially smaller than
the free diagonal term.  The finitely many remaining images are absorbed in a
constant depending only on \(T\).  Hence, for
\(\rho_\tau(y)=\int_0^\tau p_{2r}(y,y)\,dr\),
\[
        \rho_\tau(y)\asymp \sqrt{\tau}\wedge d(y),\qquad
        \int_0^1\rho_\tau(y)^{-1/2}dy\le C_T\tau^{-1/4}.      \tag{4}
\]
Indeed the integral of \(r^{-1/2}(1\wedge d^2/r)\) is comparable to
\(\sqrt\tau\) when \(d\ge\sqrt\tau\), and to \(d\) otherwise.  Also
\[
 \sup_z\int_0^1p_a(z,y)\rho_\tau(y)^{-1/2}dy
 \le C_T(a^{-1/4}+\tau^{-1/4}).                              \tag{5}
\]
To see this, use \(\rho_\tau^{-1/2}\le C(\tau^{-1/4}+d^{-1/2})\).  The first
part integrates to \(C\tau^{-1/4}\).  For the second, bound
\(d(y)^{-1/2}\le y^{-1/2}+(1-y)^{-1/2}\) and use
\(y=\sqrt a\,r\) near the left boundary (similarly near the right) to get
\(Ca^{-1/4}\sup_{\alpha\ge0}\int_0^\infty
e^{-(r-\alpha)^2/C}r^{-1/2}dr<\infty\).  Finally, for \(0<\tau\le T\),
\[
        \|\rho_\tau^{-1}S_\tau f\|_2\le C_T\tau^{-1/2}\|f\|_2. \tag{6}
\]
Here \(\rho_\tau^{-1}\le C(\tau^{-1/2}+d^{-1})\); \(S_\tau f\in H^1_0\);
Hardy's inequality gives \(\|S_\tau f/d\|_2\le C\|\partial_xS_\tau f\|_2\);
and \(\|\partial_xS_\tau f\|_2\le C\tau^{-1/2}\|f\|_2\).

Write \(D^n_{s,t}=K_t^n-S_{t-s}K_s^n\).  Freeze the equation at time \(s\):
\[
 A^n_{s,t}=\int_s^tS_{t-r}
 b_n\!\left(S_{r-s}u_s^n+V_r-S_{r-s}V_s\right)\,dr .          \tag{7}
\]
For fixed \(n\), if \(\Pi\) is a partition of \([s,t]\), then
\[
 \sum_{[a,b]\in\Pi}S_{t-b}A^n_{a,b}\to D^n_{s,t}\quad\hbox{in }L^2(H).
 \tag{8}
\]
Indeed the summand differs from
\(\int_a^bS_{b-r}b_n(u_r^n)\,dr\) by at most
\(C_n(b-a)^2\) in \(L^2(H)\), because \(b_n\) is Lipschitz and the two
arguments differ only by the drift accumulated on \([a,r]\), whose \(H\)-norm
is at most \(C_n(r-a)\).  After applying \(S_{t-b}\) and summing, the error is
\(O_n((t-s)|\Pi|)\).

The estimates below are uniform in \(n\).  Given \(\calF_s\),
\[
Z_r(y):=S_{r-s}u_s^n(y)+V_r(y)-S_{r-s}V_s(y)
\]
is Gaussian with variance \(\rho_{r-s}(y)\).  Since \(b_n\) is a probability
density, \(\sup_m\E b_n(m+N(0,\sigma^2))\le C\sigma^{-1}\).  Moreover, given
\(\calF_r\), the future part of \(Z_q(z)\), \(q>r\), has variance
\[
 \int_r^q\!\!\int_0^1p_{q-\ell}(z,\xi)^2\,d\xi d\ell
 =\int_0^{q-r}p_{2a}(z,z)\,da=\rho_{q-r}(z).
\]
Thus, for \(s<r<q\),
\[
 \E_s\,b_n(Z_r(y))b_n(Z_q(z))
 \le C\rho_{r-s}(y)^{-1/2}\rho_{q-r}(z)^{-1/2}.              \tag{9}
\]
Using
\(\int p_{t-r}(x,y)p_{t-q}(x,z)dx=p_{2t-r-q}(y,z)\), then (5) in \(y\) and
(4) in \(z\), we obtain
\[
 \E_s\|A^n_{s,t}\|_2^2
 \le C\!\int_s^t\!\!\int_r^t
      \bigl((2t-r-q)^{-1/4}+(r-s)^{-1/4}\bigr)(q-r)^{-1/4}\,dq\,dr
 \le C|t-s|^{3/2}.                                          \tag{10}
\]
The last inequality is just scaling on the triangle
\(0\le r-s\le q-s\le t-s\).  Hence
\(\|A^n_{s,t}\|_{L^2H}\le C|t-s|^{3/4}\).

For \(s<u<t\), let
\(\delta A^n_{s,u,t}=A^n_{s,t}-S_{t-u}A^n_{s,u}-A^n_{u,t}\).  The pieces before
\(u\) cancel.  For \(r>u\), the two frozen arguments differ by
\(S_{r-u}D^n_{s,u}\).  Conditional on \(\calF_u\), the remaining noise at
space point \(y\) has variance \(\rho_{r-u}(y)\), and
\(\|(b_n*\varphi_\sigma)'\|_\infty\le C\sigma^{-2}\), uniformly in \(n\).
Therefore
\[
 \left|\E_u\{b_n(Z_r(y))-b_n(Z_r(y)+h(y))\}\right|
      \le C\rho_{r-u}(y)^{-1}|h(y)|.
\]
With \(h=S_{r-u}D^n_{s,u}\), (6) and contraction of \(S_t\) imply
\[
        \|\E_u\delta A^n_{s,u,t}\|_2
        \le C(t-u)^{1/2}\|D^n_{s,u}\|_2 .                    \tag{11}
\]
Taking \(\E_s\) and using Jensen gives the same estimate with \(\E_s\) in
place of \(\E_u\).

We need only the following dyadic sewing fact.  Suppose \(A_{s,t}\) is
\(H\)-valued and \(\calF_t\)-measurable, and for every dyadic interval
\([a,b]\) with midpoint \(u\),
\[
\|\E_a\delta A_{a,u,b}\|_{L^2H}\le \Gamma_1(b-a)^{1+\varepsilon_1},\quad
\|\delta A_{a,u,b}-\E_a\delta A_{a,u,b}\|_{L^2H}
        \le \Gamma_2(b-a)^{1/2+\varepsilon_2}.
\]
If \(J_m=\sum_iS_{t-t_{i+1}}A_{t_i,t_{i+1}}\) over the \(2^m\) equal
subintervals of \([s,t]\), then \(J_m\) converges and
\[
 \|{\cal A}_{s,t}-A_{s,t}\|_{L^2H}
 \le C(\Gamma_1|t-s|^{1+\varepsilon_1}
       +\Gamma_2|t-s|^{1/2+\varepsilon_2}).                 \tag{12}
\]
Indeed, \(J_{m+1}-J_m\) is the sum of
\(-S_{t-b}\delta A_{a,u,b}\).  The conditional-mean parts are bounded by the
triangle inequality, giving a factor
\(2^m(|t-s|2^{-m})^{1+\varepsilon_1}\).  The centered parts are orthogonal:
for two level-\(m\) intervals with the first lying earlier, the earlier
summand is \(\calF_a\)-measurable for the later one, while the later summand
has conditional mean zero given \(\calF_a\).  Thus the square norm is the sum
of the square norms.  Summing the resulting geometric series proves (12).

For fixed \(n\), \(Y_n(h)=\sup_{0<t-s\le h}\|D^n_{s,t}\|_{L^2H}/|t-s|^{3/4}\)
is finite, because \(b_n\) is bounded.  On intervals of length at most \(h\),
(11) gives
\[
 \|\E_a\delta A^n_{a,u,b}\|_{L^2H}\le C Y_n(h)(b-a)^{5/4}.
\]
Also
\[
\|\delta A^n_{a,u,b}-\E_a\delta A^n_{a,u,b}\|_{L^2H}
 \le 2\|\delta A^n_{a,u,b}\|_{L^2H}
 \le C(b-a)^{3/4},
\]
because \(\delta A\) is a sum of three \(A\)-terms and (10) applies to each.
The dyadic sewing limit is \(D^n_{s,t}\) by (8).  Combining (10) and (12),
\[
 \|D^n_{s,t}\|_{L^2H}\le C|t-s|^{3/4}
        +C Y_n(h)|t-s|^{5/4}.
\]
Thus \(Y_n(h)\le C+Ch^{1/2}Y_n(h)\).  Choose \(h_0>0\), depending only on
\(T\), so that \(Ch_0^{1/2}\le1/2\); then \(Y_n(h_0)\le2C\), uniformly in
\(n\).  For longer intervals use the additivity
\(D^n_{s,t}=S_{t-r}D^n_{s,r}+D^n_{r,t}\) and split into pieces of length at
most \(h_0\).  This gives a linear bound \(C_T|t-s|\), which is bounded by
\(C_T|t-s|^{3/4}\) on \(0\le t-s\le T\).  The lemma follows.
\end{proof}

\section*{The reversible reference measure and its gap}

Let \(\gamma\) be Brownian bridge measure on \(E\), the centered Gaussian
measure with covariance \((-\Delta_D)^{-1}\).  Define
\[
  \Lambda(v)=\int_0^1\1_{\{v(r)>0\}}\,dr,\qquad
  \pi(dv)=Z^{-1}e^{2\Lambda(v)}\,\gamma(dv).                 \tag{13}
\]
For \(B_n'=b_n\), \(B_n(a)=\int_{-\infty}^ab_n(r)dr\), put
\[
   \pi_n(dv)=Z_n^{-1}\exp\left(2\int_0^1B_n(v(r))\,dr\right)\gamma(dv).
\]
Since \(0\le B_n\le1\), \(B_n(a)\to\1_{\{a>0\}}\) for \(a\ne0\), and a
Brownian bridge spends zero Lebesgue time at a fixed level,
\[
        \|\pi_n-\pi\|_{\TV}\to0 .                            \tag{14}
\]
All densities \(d\pi_n/d\gamma\) and \(d\pi/d\gamma\) are bounded above and
below by deterministic positive constants.

\begin{lemma}\label{lem:reggap}
There are constants \(C,c>0\), independent of \(n\), such that for every
bounded continuous \(f,g\) and every \(t\ge0\),
\[
 \int P_t^n f\,d\pi_n=\int f\,d\pi_n,\qquad
 \int gP_t^n f\,d\pi_n=\int fP_t^n g\,d\pi_n,                \tag{15}
\]
and
\[
        \Var_{\pi_n}(P_t^n f)\le e^{-ct}\Var_{\pi_n}(f).      \tag{16}
\]
\end{lemma}

\begin{proof}
Let \(\Pi_m\) be projection onto the first \(m\) Dirichlet modes and set
\[
        U_{n,m}(v)=\int_0^1B_n(\Pi_m v(r))\,dr .
\]
Its \(H\)-gradient is \(\nabla_HU_{n,m}(v)=\Pi_m b_n(\Pi_m v)\).  Let
\(P_t^{n,m}\) be the semigroup of
\[
u_t^{n,m}=S_t x+V_t+\int_0^tS_{t-r}\Pi_m b_n(\Pi_m u_r^{n,m})\,dr .
\]
The first \(m\) Fourier coordinates solve the finite-dimensional Langevin
diffusion
\[
dX_j=-(\lambda_j/2)X_j\,dt+\partial_j U_{n,m}(X)\,dt+d\beta_j,
\qquad j\le m,
\]
and the tail is an independent OU process.  Hence, by finite-dimensional
integration by parts and the standard OU integration by parts in the tail,
\(P_t^{n,m}\) is reversible with respect to
\(\pi_{n,m}=Z_{n,m}^{-1}e^{2U_{n,m}}\gamma\), and its Dirichlet form is
\(\calE_{n,m}(F,F)=\frac12\int\|\nabla_HF\|_2^2\,d\pi_{n,m}\) on smooth
cylinder functions.

The measures \(\pi_{n,m}\) converge to \(\pi_n\) in total variation: \(\Pi_m
v\to v\) in \(H\), \(\gamma\)-a.s., and \(B_n\) is bounded and Lipschitz.  For
fixed \(n\), \(u_t^{n,m}\to u_t^n\) in probability in \(E\).  In \(H\), this
follows from Gronwall, since \(\Pi_m b_n(\Pi_m z)\) is uniformly Lipschitz in
\(z\) for fixed \(n\), and for each \(H\)-valued \(y\),
\(\Pi_m b_n(\Pi_m y)\to b_n(y)\) in \(H\), dominated by \(2\|b_n\|_\infty\).
For convergence in \(E\) at a fixed \(t>0\), split the drift convolution into
\([0,t-\varepsilon]\) and \([t-\varepsilon,t]\).  On the first part,
\(S_{t-r}:H\to H^\alpha_0\subset E\), \(\alpha>1/2\), is bounded.  On the
short terminal part,
\[
 \int_{t-\varepsilon}^t\|S_{t-r}F_r\|_\infty dr
 \le C\|b_n\|_\infty\int_0^\varepsilon a^{-1/4}da,
\]
for both \(F_r=\Pi_m b_n(\Pi_mu_r^{n,m})\) and \(F_r=b_n(u_r^n)\), because
\(\Pi_m\) is an \(L^2\)-contraction and \(S_a:H\to E\) has norm \(Ca^{-1/4}\).
This tends to zero with \(\varepsilon\), uniformly in \(m\).

Consequently the reversible identities pass to \(m=\infty\).  For instance,
if \(h_m=d\pi_{n,m}/d\gamma\), then \(h_m\to h_n=d\pi_n/d\gamma\) in
\(L^1(\gamma)\) and the \(h_m\)'s are uniformly bounded; together with
\(P_t^{n,m}f(x)\to P_t^nf(x)\) and bounded convergence this gives (15).

It remains to record the gap uniformly.  The Gaussian Poincare inequality for
\(\gamma\) in \(H\)-gradient form is
\[
        \Var_\gamma(F)\le \lambda_1^{-1}\int\|\nabla_HF\|_2^2\,d\gamma
\]
on smooth cylinder functions.  Since all densities
\(d\pi_{n,m}/d\gamma\) lie in a fixed interval \([a_0,a_0^{-1}]\), the
Holley--Stroock bounded perturbation argument gives
\[
        \Var_{\pi_{n,m}}(F)\le C\,\calE_{n,m}(F,F),
\]
with \(C\) independent of \(n,m\).  For a symmetric Markov semigroup this
Poincare inequality is equivalent to exponential \(L^2\)-decay, so
\[
 \Var_{\pi_{n,m}}(P_t^{n,m}f)\le e^{-ct}\Var_{\pi_{n,m}}(f).
\]
Letting \(m\to\infty\), using the convergence just proved and the uniformly
bounded densities, yields (16).
\end{proof}

\begin{proposition}\label{prop:gap}
The measure \(\pi\) is \(P_t\)-invariant, \(P_t\) is symmetric on
\(L^2(\pi)\), and
\[
       \Var_\pi(P_t f)\le e^{-ct}\Var_\pi(f),
       \qquad f\in L^2(\pi),\ t\ge0 .                       \tag{17}
\]
\end{proposition}

\begin{proof}
For \(f,g\in C_b(E)\), (1), (14), (15), and the uniform boundedness of the
densities imply, by dominated convergence with respect to \(\gamma\),
\[
 \int P_tf\,d\pi=\lim_n\int P_t^nf\,d\pi_n=\int f\,d\pi,\qquad
 \int gP_tf\,d\pi=\int fP_tg\,d\pi .
\]
On the Polish space \(E\), these identities identify the measure
\(\int\pi(dx)P_t(x,\cdot)\) with \(\pi\), and the measure
\(\pi(dx)P_t(x,dy)\) with its transpose; hence invariance and detailed balance
hold for all bounded Borel functions.  Similarly, (16) passes to the limit for
bounded continuous \(f\).  Finally \(C_b(E)\) is dense in \(L^2(\pi)\), and an
invariant Markov operator is an \(L^2\)-contraction, so symmetry and (17)
extend to all \(L^2(\pi)\).
\end{proof}

\section*{Fixed-time absolute continuity}

\begin{lemma}\label{lem:control}
For every deterministic \(x\in E\) and every \(t>0\),
\[
        P_t(x,\cdot)\ll \gamma ,
\]
and hence \(P_t(x,\cdot)\ll\pi\).
\end{lemma}

\begin{proof}
Let \(\nu_t^x=\Law(S_tx+V_t)\).  We first bound
\(\Ent(P_t^n(x,\cdot)\mid\nu_t^x)\) uniformly in \(n\).  Put
\(r_k=t(1-2^{-k})\), \(\Delta_k=r_{k+1}-r_k\), and
\(D_k^n=K^n_{r_{k+1}}-S_{\Delta_k}K^n_{r_k}\).  Define
\[
 F_s^n=\sum_{k\ge0}\frac{\1_{(r_{k+1},r_{k+2}]}(s)}{\Delta_{k+1}}\,
        S_{s-r_{k+1}}D_k^n .                                \tag{18}
\]
The process is predictable because \(D_k^n\) is \(\calF_{r_{k+1}}\)-measurable.
Lemma \ref{lem:Hestimate} and contraction of \(S_t\) give
\[
 \E\int_0^t\|F_s^n\|_2^2ds
 \le C_t\sum_{k\ge0}\frac{\Delta_k^{3/2}}{\Delta_{k+1}}<\infty, \tag{19}
\]
uniformly in \(n\).  For partial sums,
\[
 \int_0^tS_{t-s}F_s^{n,N}ds
 =\sum_{k=0}^NS_{t-r_{k+1}}D_k^n
 =S_{t-r_{N+1}}K^n_{r_{N+1}},
\]
where \(F^{n,N}\) denotes the sum in (18) up to \(k=N\).  Since
\(r_{N+1}\uparrow t\), the regularized mild solution gives \(K^n\) continuous
with values in \(E\), and \(S_a\to I\) strongly on \(E\),
\[
        \int_0^tS_{t-s}F_s^n\,ds=K_t^n                       \tag{20}
\]
in probability in \(E\) (and hence in \(H\)).  The state space assertion is
legitimate because, for every deterministic \(F\in L^2(0,t;H)\),
\[
 \sum_j\lambda_j^\alpha\Big|\int_0^te^{-\lambda_j(t-s)/2}F_j(s)ds\Big|^2
 \le \sum_j\lambda_j^{\alpha-1}\int_0^t|F_j(s)|^2ds
 \le C\|F\|_{L^2(0,t;H)}^2
\]
for \(0<\alpha<1\); choosing \(\alpha>1/2\) gives an element of
\(H^\alpha_0(0,1)\subset E\).

We use the following standard entropy form of Girsanov, and recall its proof.
If \(F\) is predictable, \(H\)-valued, and has bounded energy, then
\[
 X^F=S_tx+V_t+\int_0^tS_{t-s}F_sds
\]
satisfies
\[
        \Ent(\Law(X^F)\mid\nu_t^x)
        \le \frac12\,\E\int_0^t\|F_s\|_2^2ds .               \tag{21}
\]
For the first \(M\) Brownian coordinates this is finite-dimensional
Girsanov: under the tilted measure, the shifted coordinates
\(\beta_j+\int_0^\cdot F_j(s)ds\), \(j\le M\), are Brownian, and the entropy
cost is \(\frac12\E\int_0^t\sum_{j\le M}F_j(s)^2ds\).  Applying data
processing to the deterministic-kernel stochastic convolution gives the same
bound for the first \(M\) terminal Fourier coordinates.  Since the Borel
sigma-field of \(E\) is generated by these coordinates and relative entropy
is the supremum over increasing finite-coordinate sigma-fields, (21) follows.

Localize \(F^n\) by
\(\tau_R=\inf\{s:\int_0^s\|F_r^n\|_2^2dr>R\}\) and
\(\widehat F^{\,n,R}=F^n\1_{[0,\tau_R]}\).  Let
\[
 X_R=S_tx+V_t+\int_0^tS_{t-s}\widehat F^{\,n,R}_sds .
\]
Then (21) and (19) give
\[
 \Ent(\Law(X_R)\mid\nu_t^x)
 \le \frac12\,\E\int_0^t\|F_s^n\|_2^2ds .                   \tag{22}
\]
On \(\{\tau_R>t\}\), (20) gives \(X_R=u_t^n\), while
\(\mathbb P(\tau_R\le t)\le R^{-1}\E\int_0^t\|F_s^n\|_2^2ds\to0\).  Hence
\(X_R\to u_t^n\) in probability in \(E\).  Lower semicontinuity of relative
entropy on the Polish space \(E\), then (19) and (22), yield
\[
        \sup_n\Ent(P_t^n(x,\cdot)\mid\nu_t^x)<\infty .       \tag{23}
\]
By the assumed weak convergence (1) and lower semicontinuity again,
\(\Ent(P_t(x,\cdot)\mid\nu_t^x)<\infty\).  Thus
\(P_t(x,\cdot)\ll\nu_t^x\).

It remains only to compare Gaussian measures.  In the basis
\(e_j(r)=\sqrt2\sin(j\pi r)\), \(\lambda_j=j^2\pi^2\), the covariance of
\(\nu_t^x\) has eigenvalues
\[
        q_j(t)=\frac{1-e^{-\lambda_jt}}{\lambda_j},
\]
whereas \(\gamma\) has eigenvalues \(\lambda_j^{-1}\).  The ratios
\(1-e^{-\lambda_jt}\) are bounded away from zero and differ from \(1\) by a
square-summable sequence.  The Cameron--Martin spaces are therefore both
\(H^1_0(0,1)\), and \(S_tx\in H^1_0\).  Feldman--Hajek gives
\(\nu_t^x\sim\gamma\).
\end{proof}

\section*{Conclusion}

Let \(\eta\) be an invariant probability measure.  If \(\pi(A)=0\), then
Lemma \ref{lem:control} gives \(P_t(y,A)=0\) for every \(y\in E\), and hence
\[
        \eta(A)=\int_E P_t(y,A)\,\eta(dy)=0.
\]
Thus \(\eta\ll\pi\); write \(d\eta=h\,d\pi\).  Detailed balance gives the
\(L^1\)-duality
\[
        \int f\,P_t g\,d\pi=\int P_t f\,g\,d\pi
\]
for bounded Borel \(f\) and \(g\in L^1(\pi)\), by truncating \(g\).  Therefore
\[
 \int f h\,d\pi=\int P_tf\,d\eta
      =\int P_tf\,h\,d\pi=\int f\,P_t h\,d\pi ,
\]
so \(P_t h=h\) in \(L^1(\pi)\).  For \(a>0\), let \(h_a=h\wedge a\).  Since
\(r\mapsto r\wedge a\) is concave,
\[
        P_t h_a\le (P_t h)\wedge a=h_a .
\]
The two sides have the same \(\pi\)-integral; therefore \(P_t h_a=h_a\).
Now \(h_a\in L^2(\pi)\), and the spectral gap (17) implies
\(\Var_\pi(h_a)=0\).  Thus \(h_a\) is constant for every \(a\).  Letting
\(a\uparrow\infty\) and using \(\int h\,d\pi=1\), we obtain \(h=1\)
\(\pi\)-a.s.  Hence \(\eta=\pi\).

Therefore \(P_t\) has at most one invariant probability measure on
\(C_0([0,1])\).  In fact, if an invariant probability exists, it is the Gibbs
measure \(\pi\) in (13).

\begin{thebibliography}{9}

\bibitem{ABLM}
S. Athreya, O. Butkovsky, K. L\^e and L. Mytnik,
\emph{Well-posedness of stochastic heat equation with distributional drift
and skew stochastic heat equation}, Comm. Pure Appl. Math. 77 (2024),
2708--2777.

\bibitem{LeBanach}
K. L\^e,
\emph{Stochastic sewing in Banach spaces}, Electron. J. Probab. 28 (2023),
paper no. 26, 1--22.

\bibitem{LeSSL}
K. L\^e,
\emph{A stochastic sewing lemma and applications}, Electron. J. Probab. 25
(2020), paper no. 38, 1--55.

\bibitem{Bogachev}
V. I. Bogachev,
\emph{Gaussian Measures}, American Mathematical Society, 1998.

\bibitem{DPZ}
G. Da Prato and J. Zabczyk,
\emph{Ergodicity for Infinite Dimensional Systems}, Cambridge University
Press, 1996.

\bibitem{MR}
Z.-M. Ma and M. R\"ockner,
\emph{Introduction to the Theory of (Non-Symmetric) Dirichlet Forms},
Springer, 1992.

\bibitem{BZ}
S. K. Bounebache and L. Zambotti,
\emph{A skew stochastic heat equation}, J. Theoret. Probab. 27 (2014),
168--201.

\end{thebibliography}

\end{document}
